%Keywords: differential equations, stationary point, resonant terms,
%Hamiltonian system.

\documentstyle[12pt]{article}
\setlength{\evensidemargin}{0in}
\setlength{\oddsidemargin}{0in}
\setlength{\textwidth}{16.cm}
\setlength{\textheight}{8.5in}
\setlength{\topmargin}{0in}
\setlength{\headheight}{0in}
\setlength{\headsep}{0in}
\setlength{\itemsep}{-\parsep}
\renewcommand{\topfraction}{.9}
\renewcommand{\textfraction}{.1}
\newcommand{\ol}{\setlength{\itemsep}{0pt.}\begin{enumerate}}
\newcommand{\eol}{\end{enumerete}\setlength{\itemsep}{-\parsep}}
\newcommand{\R}{{\rm{I\! R}}}
\newcommand{\Q}{{\rm Q\!\!\! l}}
\newcommand{\C}{{\rm{I\!\!\! C}}}
\newcommand{\N}{{\rm{I\! N}}}
\newcommand{\defi}{{\;\buildrel \rm def \over =}\;}
\hyphenation{tra-ns-for-ma-tion}

\begin{document}
\Large
\centerline {\bf On computation of normal forms}
\centerline {\bf A.D.~Bruno}
\centerline {Institute of Applied Mathematics, Moscow}
\centerline { 125047 Russia (bruno@applmat.msk.su)}
\bigskip
\bigskip

{\bf Keywords:} differential equations, stationary point, resonant terms,
Hamiltonian system.

Let $X\in \R^n$ or $\,\C^n$. In a neighbourhood of the stationary point $X=0$ 
we
consider the ODE system
$$
dX/dt\defi \dot X=AX+\Phi(X)\eqno(1)
$$
where $\Phi$ is a vector power series, $\Phi =O(\|X\|^2)$.

{\bf Problem:} Simplify the system (1) preserving  main properties of
its solutions.

There are two methods for that:

1. If in (1) the matrix $A\neq 0$, then we can perform a formal
transformation of coordinates
$$
X=BY+ \Xi(Y)\eqno(2)
$$
to bring the system (1) into a normal form
$$
\dot Y = CY+\Psi(Y).\eqno(3)
$$

2. If $A=0$ then instead of the system (1) we can consider  the
truncated system
$$
\dot X = \hat\Phi_j^{(d)}(X),\eqno(4)
$$
where components of the vector $\hat\Phi_j^{(d)}(X)$ are some
quasihomgeneous parts of components of the vector-series $\Phi(X)$.
Generally speaking, there are several different truncated systems (4).
To find them we use the Newton polyhedron of the system (1).

In my book [1] it was shown that in any compicated case,
combining both methods, one can reduce the system (1) to several
simple systems. Here we shall only consider the first method in
details.

Let the matrix $C=(c_{ij})$ be in a normal form (for instance, in the
Jordan one). So its diagonal consists of its eigenvalues
$(\lambda_1,\dots,\lambda_n)=\Lambda$, and $c_{ij}=0$ if
$\lambda_i\neq\lambda_j$. Let us write components of the vector $\Psi$
in (3) in the form
$$
\psi_i(Y)\defi y_ig_i(Y)\defi  y_i\sum g_{iQ}Y^Q,\quad i=1,\dots,n,\eqno(5)
$$
where $Q=(q_1,\dots,q_n)\in {\bf Z}^n$ and $Y^Q=y_1^{q_1} \dots y_n^{q_n}$.

{\bf Definition 1} [2,1]. The system (3), (5) is called {\it the
resonant normal form}, if expansions (5) contain only such terms
$g_{iQ}Y^Q$ for which the scalar product
$$
\langle Q,\Lambda\rangle \defi
q_1\lambda_1+\dots+q_n\lambda_n=0.\eqno(6)
$$

{\bf Theorem 1} [2,1]. {\it For each system (1) there exists a formal
transformation (2)  which brings (1) into its resonant normal form
(3).}

Let the equation (6) have exactly $k$ integer solutions $Q\in {\bf
Z}^n$ which  are linearly independent over $\R$. Then the system (1)
has the $k$ -- {\it fold resonance.}

{\bf Definition 2} [3,1]. The nonlinear transformation of coordinates
$$
z_i = y_1^{\alpha_{i1}}y_2^{\alpha_{i2}} \dots y_n^{\alpha_{in}}, \quad i=1, 
\dots, n
\eqno(7)
$$
is called {\it the power transformation} with the matrix $\alpha
=(\alpha_{ij})$. It is a linear transformation of logarithms of
coordinates $\log Z = \alpha\log Y$ where $\log Z = (\log z_1,\dots, \log 
z_n)^*$.

{\bf Theorem 2} [2,1]. {\it If the system (1)
has the $k$ --  fold resonance, then there exists such a power
transformation (7) which brings the resonant normal form (3) into the
system}
$$
\dot{(\log z_i)}=f_i(z_1,\dots,z_k),\quad i=1,\dots,n.\eqno(8)
$$
{\it Here all $\alpha_{ij}$ are integers and}  det$\alpha =\pm 1$.

Thus, to solve the system (8) one must solve its subsystem of order $k$:
$$
\dot z_i = z_if_i(z_1,\dots,z_k),\quad i=1,\dots,k.\eqno(9)
$$
The system (9) has not a linear part.
So to simplify the system (9) we must find its truncated systems
$$
\dot z_i = z_i\hat f_{ij}^{(d)}(z_1,\dots,z_k),\quad i=1,\dots, k,\eqno(10)
$$
using the Newton polyhedron of the system (9). The system (10) is
quasihomogeneous, so we can again use a power transformation to  simplity it.

There are  3 other kinds of the normal form (shortly n.f.):

2. Belitski n.f.

3. Hamiltonian resonant n.f.

4. Hamiltonian Belitski n.f.

{\bf Definition 3} [4,5]. The system
$$
\dot Y =CY+\tilde\Psi(Y)\eqno(11)
$$
is called the {\it Belitski normal form}, if it is the resonant n. f.
and
$$
{\partial\tilde \Psi\over\partial Y}\bar{C^*}Y-\bar{C^*}\tilde \Psi(Y) =
0,\eqno(12)
$$
where the matrix $\bar{C^*}$ is conjugate to the matrix $C$.

If components of the vector $\tilde \Psi$ we write  as
$$
\tilde \psi_i = y_i\sum\tilde g_{iQ}Y^Q,\quad i=1,\dots, n
$$
then these expansions for the Belitski n. f. (11)
have only resonant terms $\tilde g_{iQ}Y^Q$ with property (6) but some
of them are absent.

{\bf Theorem 3} [4,5]. {\it For each system (1) there exists a formal
transformation (2) into its Belitski n. f.}

The power transformation (7) of Theorem 2 brings the Belitski n. f.
(11) into the system
$$
(\log \dot z_i) =\tilde f_i(z_1,\dots, z_k),\quad i=1,\dots,n
$$
with the subsystem
$$
\dot z_i = z_i\tilde f_i(z_1,\dots, z_k),\quad i=1,\dots, k.
$$

Let us compare its truncated systems
$$
\dot z_i = z_i\widehat{\tilde f}_{ij}^{(d)}(z_1,\dots, z_k),\quad
i=1,\dots, k.  \eqno(13) $$
with truncations (10) of the system (9), obtained from the resonant
n. f. (3). Some of them can conside:
$$
\widehat f_{ij}^{(d)}(Y)=
\widehat{\tilde f}_{ij}^{(d)}(Y),  \quad i=1,\dots, k.
 $$
 Thus, to compute such truncations it is enough to compute the
 resonant n. f. that is simpler than a computation of the Belitski
 n.f.

 {\bf Example 1.} In [6] there was considered an invertible system (1)
 with  $n=4$, where the matrix $A$ depends of two small parameters:
 $A=A(\mu_1, \mu_2)$. The unperturbed matrix $A=A(0,0)$ is the
 $4\times 4$ Jordan block  with zero eigenvalue. Then the system (1)
 is its own resonant n.f. By means of a special program [7]  for
 computation of the Newton polyhedron we have computed the Newton
 polyhedron of that system (1). It gave 5 truncated systems; 4 of them
 are integrable and one (the main  truncation) is
 nonintegrablite. We have computed   the Belitski n.~f.
 of the system (1) and its Newton polyhedron. It gave 4 truncated
 systems for the Belitski n. f.; 3 of them are integrable and the
 fourth is exactly the same main truncation that was obtained for
 initial system (1).

 Now let us consider the Hamiltonian system (1):
 $$
 \dot x_i  = {\partial h\over\partial x_{i+m}},\quad
 \dot x_{i+m}  = -{\partial h\over\partial x_{i}},\quad
 i=1,\dots, m\eqno(14)
 $$
 where $n=2m$ and the Hamiltonian function $h$ is a power series
 $$
 h=\sum h_QX^Q.\eqno(15)
 $$
 Let $\Lambda$ be the vector of eigenvalues of the matrix of the
 linear part of the system (14).

 {\bf Definition 4 }[8]. {\it Hamiltonian} system (14) is called the
 {\it resonant normal form} if the expansion (15) contains only terms
 $h_QX^Q$ with the property (6).

 {\bf Theorem 4} [8, \S~12]. {\it For any Hamiltonian system (14), there
 exists a canonical transformation (2) which brings (14) into the
 Hamiltonian resonant n. f.}

 For the Hamiltonian resonant n. f., one can compute the reduced
 subsystem
 $$
 \dot z_i = z_i\tilde h_i(z_1,\dots, z_k),\quad i=1,\dots, k.\eqno(16)
 $$
 On the other side, for the Hamiltonian system (14), we can compute
 the resonant n. f. (3), using a noncanonical transformation (2) and
 obtain its reduced sybsystem  (10). Sometimes the main truncations of
 the systems (10) and (16) differ by a linear transformation of
 variables only. In such cases, to obtain the
 main truncation we can use the
 noncanonical normalizing
 transformation, that demands essentially more simple computation that
 the canonical one. Moreover, to compute coefficients of the
 truncated system (10) we can use the explicit formula (35) in [1, Ch.
 III, \S~1]. In [9] such approach was used in a system related to  the
 water--wade problem.

 Let $C$ be the matrix of the linear part of the Hamiltonian system
 (14).

 {\bf Definition 5 }[5]. The Hamiltonian system (14) is called {\it
 the Hamiltonian Belitski normal form} if the scalar product
 $$
 \langle \bar C^*X, \partial h/\partial X\rangle = 0.
 $$

 {\bf Theorem  5 }[5]. {\it There exists a formal canonical transformation
 (2), which brings the Hamiltonian system (14) into its Hamiltonian
 Belitski n. f.}

 Indeed the normal form and the normalizing transformation are not
 unique. Using more special normalizing transformation, we obtain more
 special normal form, but it demands more complicated computation.

 Thus,  me have 4 normal forms:

 (i)  Resonant n. f.

 (ii) Belitski n. f.

 (iii) Hamiltonian resonant n. f.

 (iv) Hamiltonian Belitski n. f.

 There are inclusions:
 $$
 \begin{array}{c}
 ({\rm i}) \supset ({\rm iii})\supset({\rm iv}),\\
 ({\rm i}) \supset ({\rm ii})\supset({\rm iv}).
 \end{array}
 $$

 Apparently in some cases we can solve a problem, computing less
 special n. f., i.e. using more simple computation. Five methods of
 computation of normal forms are discussed in [10].

 The work was partly supported by the  International Science
 Foundation, Grant M3W300.


\begin{thebibliography}{99}
\bibitem{1}
    Bruno, A.D. Local Methods in Nonlinear Differential Equations.
   Springer -Verlag, Berlin, 1989.
\bibitem{2}
   Bruno,  A.D. Normal form of differential equations.
   Soviet Math. Dokl. {\bf 5}(1964) 1105 -- 1108.
\bibitem{3}
   Bruno,  A.D. The asymptotic behavior of solutions of nonlinear
   systems of differential equations.
   Soviet Math. Dokl. {\bf 3}(1962) 464 -- 467.
\bibitem{4}
    Belitski, G.R. The normal forms of local mappings. Uspehi Mat.
    Nauk, {\bf 30}:1(1975) 223 (Russian).
\bibitem{5}
    Belitski, G.R. Normal Forms, Invariants and Local Mappings.
    Kiev: Naukova Dumka, 1979 (Russian).

\bibitem{6}
 Soleev, A. and Aranson, A.B. First approximations of a reversible ODE
  system. Preprint N~28 of the Keldysh Inst. of Appl. Math., Moscow, 1995
 (Russian).
\bibitem{7}
 Soleev, A. and Aranson, A.B. Computation of a polyhedron and of
 normal cones of its faces. Preprint N~36 of the Keldysh Inst. of
 Appl. Math., Moscow, 1994 (Russian).

\bibitem{8}
   Bruno, A.D. Analytical form of differential equations. Trans. Moscow
   Math. Soc.{\bf 26}(1972) 199--239.

\bibitem{9}
  Bruno, A.D. and  Soleev, A.  Local analysis of a singularity of an
  invertible ODE system. Preprints Nos. 40, 47, 54 of Keldysh Inst. of
  Appl. Math., Moscow, 1995 (Russian).

\bibitem{10}
   Bruno, A.D. Methods of computation of the normal form. Russian
   Academy Sci. Doklady. Mathem. (1995) (to appear).
\end{thebibliography}

\end{document}
